\section{Related work}
\label{sec:related}

\paragraph{Failure and inefficiency in agent trajectories.}
Recent work has moved from final success rates to process-level diagnostics. An anatomy study of CLI
coding-agent trajectories annotated thousands of runs and found that failures are predominantly
epistemic, usually begin within the first few steps, and often stay hidden until recovery is no longer
possible \cite{zhao2026failure}. Trace-structure diagnostics argue that resolve rate hides behavioural
pathologies, naming anti-patterns such as search loops and verification skips \cite{shu2026traceprobe}. We
share that diagnostic motivation but differ in what we compute: rather than classifying trajectories after
the fact, we estimate at every step whether the last window of activity moved the task, using signals a
runtime can compute from the prefix alone.

\paragraph{Early failure prediction and early termination.}
Fail-Fast/Restart-Smart trains a small monitor on terminal and dense fail-to-pass supervision and restarts
the trajectory when it alarms, reporting token savings at a target false-positive rate
\cite{wang2026failfast}; AgentStop terminates local agents on platform-level signals to save energy
\cite{pham2026agentstop}. Both are close in spirit to our alarm setting, and we adopt their central
evaluation device --- savings at a fixed false-stop rate. Two differences matter. Their supervision is the
final outcome, so their monitor learns ``this run will fail'', whereas our labels are independent window
annotations of whether progress happened, and Section~\ref{sec:results} shows the two are not
interchangeable. And our features are symbolic and reference-free, which makes each alarm explainable by
the entities and verification deltas that produced it.

\paragraph{Livelock and semantic redundancy.}
The closest antecedent frames the phenomenon as \emph{semantic livelock} --- an agent that keeps generating
tokens and calling tools without progressing --- and reports that it accounts for $25\%$ of the
long-duration failures in a forensic analysis of SWE-agent runs, including one agent that spent $208$ steps
in a checkerboard oscillation invisible to string-matching repetition guards \cite{khanna2026zombie}. That
work proposes a sidecar that fingerprints agent state in embedding space. Our corpus and measurements
differ in kind, so the comparison is specific rather than general: they characterise how often livelock
occurs and propose a monitor, whereas we build a labelled benchmark and ask \emph{which} observable signal
detects it. Two of our results bear on that proposal. Embedding-space similarity between successive steps
is the strongest single family we tested at window level (\paucBFoursemantic{} ROC-AUC), which independently
supports the choice of representation --- but it is not the signal we could alarm on, because the semantic
score rises monotonically with run progress and no fixed threshold on it stayed calibrated; a sidecar built
on embedding similarity therefore needs an explicit treatment of drift that the proposal does not address.
We also quantify what the string-matching guard misses and find it narrower than expected: exact repetition
still reaches \paucBTworepThree{}, so the gap between naive guards and semantic ones is real but not the
order of magnitude the oscillation example suggests.

\paragraph{Credit assignment, verification, benchmarks and data.}
Transition-wise credit assignment evaluates each transition with evidence, execution and invalidity
rubrics, and rewards newly covered conditions \cite{zhang2026trca}; verification frameworks score candidate
solutions and track progress toward them \cite{kwok2026verifier}. These are training-time or offline
instruments, and our evidence channel is the cheap deterministic cousin of their evidence rubric: instead
of asking a model whether a transition acquired task-relevant information, we check whether an entity the
observation exposed is new and lexically connected to the task statement. SWE-bench established
repository-level issue resolution as an evaluation target \cite{jimenez2024swebench}, and SWE-agent showed
how much the agent--computer interface matters \cite{yang2024sweagent}; we use Terminal-Bench~2.0
trajectories \cite{yoonholee2026tb2} because they cover many scaffolds and models on a common task set, and
SWE-agent trajectories \cite{nebius2024sweagenttraj} for a cross-corpus check. Production issue trackers
document the practical failure mode this paper targets: agents that keep investigating without progress,
long runs of identical calls with no loop protection, and stall watchdogs that kill healthy requests
because they cannot tell stagnation from slow work \cite{opencode43603,opencode45442}.

\paragraph{Position of this paper.}
We propose no new learning method. We take the signals the literature suggests --- repetition, semantic
diversity, novelty, verification movement, credit-style evidence --- implement them all on the same feature
matrix, and measure them against the same independent labels under the same constraint (online,
reference-free, task-disjoint training). The contribution is the comparison and the negative results, plus
a task-grounded evidence signal cheap enough to run at every step.
